\documentclass{situationist-preprint}

\title{On the Indiscriminate Posing of Problems}
\msid{EPI-22-b}
\msfamily{EPI-22}
\msdiscipline{Epistemology}
\mstype{Consequence}
\mskeywords{curiosity; intellectual virtue; problem-posing; significance; attention; inquiry}

\begin{document}
\maketitle

\begin{abstract}
EPI-22-a argues that curiosity is an epistemic virtue when it is directed at significant questions, and that indiscriminate inquiry is not virtuous. We derive a consequence for a related practice: the posing of problems to others. To pose a problem is to direct other inquirers' finite attention to a question, with an implicit claim that it deserves it. We show that problem-posing that is indifferent to significance is not virtuous, for two reasons: it manifests indiscriminate curiosity at one remove, and it imposes on others the cost in significant knowledge foregone that EPI-22-a identifies for the indiscriminate inquirer. We show further that indiscriminate posing degrades the testimonial value of the claim of significance that posing carries. Virtuous problem-posing selects by significance and makes the case for it, as the best-known lists of problems have done.
\end{abstract}

\section{Posing as directing attention}

Inquiry begins with questions, and many questions are first raised by someone other than the person who pursues them. Teachers set problems, journals solicit them, and some mathematicians and scientists are remembered as much for the problems they posed as for those they solved. The epistemology of curiosity has concentrated on the inquirer's desire to know \citep{Zagzebski1996,Whitcomb2010}. We consider the poser.

\begin{definition}[Posing a problem]\label{d:pose}
To \emph{pose a problem} is to formulate a question publicly and present it to other inquirers as worth their pursuit.
\end{definition}

\begin{assumption}[Uptake]\label{a:uptake}
A problem that is posed is, with some probability, taken up: some inquirers to whom it is presented devote attention to it that they would otherwise have devoted to other questions.
\end{assumption}

Assumption~\ref{a:uptake} combines the observation that posed problems are often pursued with Assumption~2.4 of EPI-22-a, that attention is finite: inquiry into one question is foregone inquiry into others. It follows that posing is not costless. Each problem taken up displaces some other inquiry, and the epistemic value of the displaced inquiry depends, by Principle~P1 of EPI-22-a, on the significance of the questions it would have answered.

\begin{definition}[Indiscriminate posing]\label{d:indisc}
A disposition to pose problems is \emph{indiscriminate} if which questions it leads the poser to present as worth pursuit does not depend on their significance (Definition~2.3 of EPI-22-a).
\end{definition}

\section{The consequence}

\begin{theorem}[Indiscriminate posing is not virtuous]\label{t:pose}
A disposition to pose problems indiscriminately is not an intellectual virtue.
\end{theorem}

\begin{proof}[Argument]
By Definition~2.2 of EPI-22-a, an intellectual virtue requires (i) a concern for epistemic goods proportioned to their value and (ii) reliable success in attaining them. The poser's characteristic concern, in posing, is that others come to know the answers to the questions posed. A disposition that presents questions as worth pursuit independently of their significance manifests a concern for the answers that is not proportioned to their value, since by Principle~P1 of EPI-22-a that value increases with significance. Condition (i) fails, as it does for indiscriminate curiosity in Theorem~3.2 of EPI-22-a; the disposition is indiscriminate curiosity exercised at one remove, through others' inquiry. Condition (ii) fails in expectation: by Assumption~\ref{a:uptake}, the questions taken up displace others, and when the posed questions are not selected by significance, the expected significance of what is displaced is not less than that of what is posed. The disposition therefore does not reliably increase, and may decrease, the epistemic value of the inquiry it directs.
\end{proof}

The cost identified in the second half of the argument falls not on the poser but on those who take up the problems. That makes indiscriminate posing worse, in one respect, than indiscriminate curiosity: the indiscriminately curious inquirer squanders her own attention, while the indiscriminate poser spends others'.

\begin{proposition}[Degraded signal]\label{p:signal}
Presenting a question as worth pursuit carries an implicit claim of significance. If a poser presents questions indiscriminately, that claim, made by him, ceases to be evidence of significance.
\end{proposition}

\begin{proof}[Argument]
A claim is evidence for what it says to the degree that its being made covaries with the truth of what it says. If which questions a poser presents does not depend on their significance (Definition~\ref{d:indisc}), then his presenting a question does not covary with its significance, and his implicit claim carries no evidence of it.
\end{proof}

Proposition~\ref{p:signal} has a social consequence. Inquirers rely on posers' claims of significance to allocate attention, since they cannot assess the significance of every question themselves. The indiscriminate poser deprives them of this guidance for his own problems and, where his problems circulate among others', dilutes the guidance that others' selections provide.

\section{Virtuous posing}

\begin{corollary}[Virtuous posing]\label{c:virtuous}
A disposition to pose problems is a candidate for intellectual virtue only if the questions it presents are selected by their significance; and its exercise is better when the poser makes the case for that significance, so that others can assess the claim rather than only trust it.
\end{corollary}

The corollary describes the best-known practice. Hilbert's address of 1900 posed twenty-three problems, and opened by asking what makes a problem fruitful: it should be clear, difficult enough to attract and not so difficult as to mock effort, and connected with the development of the field \citep{Hilbert1902}. Whatever one thinks of particular selections, the address presented its problems with reasons for their significance, and its influence on twentieth-century mathematics illustrates the value that well-selected problems can have. Kitcher's account of significance, on which the significance of scientific questions depends on their connection with other questions and with human interests, gives the same practice a general rationale \citep{Kitcher2001}.

\section{Objections}

\begin{objection}
Significance is often unpredictable. Many questions that seemed idle turned out to matter, and a poser who confines himself to questions of evident significance will miss them.
\end{objection}

\begin{reply}
EPI-22-a answers the parallel objection for curiosity: calibrated exploration of questions whose significance is uncertain is itself directed at expected significance. The same holds of posing. A poser may present a question as worth a modest exploratory effort, stating that its significance is uncertain; that is a discriminating act. What Theorem~\ref{t:pose} excludes is posing whose selection does not depend on significance, expected or actual, at all.
\end{reply}

\begin{objection}
Inquirers are free to ignore problems that do not interest them, so the poser imposes nothing.
\end{objection}

\begin{reply}
Ignoring a problem requires assessing it, and assessment consumes attention. Where problems are many and assessments costly, inquirers fall back on the posers' implicit claims of significance; by Proposition~\ref{p:signal}, indiscriminate posing makes those claims worthless. Freedom to ignore reduces the cost of indiscriminate posing without removing it.
\end{reply}

\begin{objection}
The tradition from Augustine onwards condemned curiosity as a vice of attention to what does not concern us \citep[X.35]{AugustineConf}. Theorem~\ref{t:pose} is that condemnation applied to posers, and inherits its moralism.
\end{objection}

\begin{reply}
The condemnation targeted curiosity as such; EPI-22-a and the present argument condemn only the indiscriminate form, and the criterion is epistemic significance, not moral or religious propriety. Theorem~\ref{t:pose} leaves room, as the Augustinian tradition did not, for the virtuous poser of questions that matter.
\end{reply}

\section{Conclusion}

Posing a problem directs other inquirers' finite attention and claims that the problem deserves it. It is therefore subject to the same constraint as curiosity: directed at significant questions it can be virtuous, and indiscriminate it is not. The indiscriminate poser spends others' attention without regard to what it buys and empties his own claims of significance of evidential force. The virtuous poser selects, and says why.

\begin{thebibliography}{99}\small
\bibitem[Augustine(Conf.)]{AugustineConf} Augustine. \emph{Confessions}. Cited by book and chapter.
\bibitem[Hilbert(1902)]{Hilbert1902} Hilbert, D. (1902). Mathematical problems. Trans. M.~W. Newson. \emph{Bulletin of the American Mathematical Society}, 8(10), 437--479.
\bibitem[Kitcher(2001)]{Kitcher2001} Kitcher, P. (2001). \emph{Science, Truth, and Democracy}. Oxford University Press.
\bibitem[Whitcomb(2010)]{Whitcomb2010} Whitcomb, D. (2010). Curiosity was framed. \emph{Philosophy and Phenomenological Research}, 81(3), 664--687.
\bibitem[Zagzebski(1996)]{Zagzebski1996} Zagzebski, L.~T. (1996). \emph{Virtues of the Mind}. Cambridge University Press.
\end{thebibliography}

\end{document}
